% !TEX root = ../../main.tex
% C5.5 — Thiết kế bảo mật và phân quyền (RÚT GỌN 2026-09-29: bảng giới hạn tần suất -> câu văn; gộp 5.5.4 + 5.5.5; bản cũ ở report/archive/)
% Khớp: auth/tokens.py (mã phiên 32 byte ngẫu nhiên, lưu SHA-256; CSRF = HMAC-SHA256(mã phiên, ngữ cảnh)); auth/web.py (cookie HttpOnly, SameSite=Lax,
% Secure ở môi trường chính thức; tiêu đề X-CSRF-Token cho phương thức không an toàn; require_auth(vai trò) -> 403);
% auth/passwords.py (Argon2, dài tối thiểu 8, băm giả khi không có tài khoản); services/auth.py (login thu hồi phiên cũ; mỗi yêu cầu kiểm tra thu hồi,
% hạn 12 h, 30 phút không hoạt động, tài khoản ACTIVE, Operator có khu vực; refresh cấp mã mới); api/security.py (CORS danh sách cho phép, tiêu đề bảo mật);
% api/rate_limit.py + config.py (login 10 theo IP + tên, search 30, upload 10, connection_test 10, diagnostics 6 mỗi phút; bộ đếm trong bộ nhớ);
% services/cases.py (không phải chủ -> 404); services/camera_runtime.py (Fernet, chỉ IP literal thuộc dải cho phép, chặn loopback/link-local/multicast).
% LƯU Ý: refresh() đặt hạn mới = now + 12 h -> hạn 12 h không còn "tuyệt đối" khi người dùng hoạt động liên tục (architect.md §13 #8). Văn bản mô tả đúng code.
\section{Thiết kế bảo mật và phân quyền}
\label{sec:5.5}

Các yêu cầu bảo mật ở Mục~\ref{sec:4.3.1} và ma trận phân quyền ở Mục~\ref{sec:4.6} được hiện thực theo nguyên tắc \emph{phòng thủ nhiều lớp}: mỗi yêu cầu được kiểm tra phiên, vai trò và phạm vi dữ liệu, và ràng buộc cơ sở dữ liệu là lớp cuối cùng.

\subsection{Xác thực và phiên đăng nhập}
\label{sec:5.5.1}

Mật khẩu dài tối thiểu 8 ký tự và được băm bằng Argon2. Tên đăng nhập không tồn tại vẫn được kiểm tra trên một giá trị băm giả và nhận cùng thông báo như sai mật khẩu, nên thời gian phản hồi và nội dung lỗi không cho biết tên nào tồn tại.

Đăng nhập thành công tạo một mã phiên ngẫu nhiên 32 byte trong cookie \texttt{HttpOnly} (mã JavaScript không đọc được), \texttt{SameSite=Lax} (không gửi kèm yêu cầu thay đổi dữ liệu từ trang khác) và \texttt{Secure} khi chạy qua HTTPS. Cơ sở dữ liệu chỉ lưu mã băm SHA-256 của mã phiên, nên bảng phiên bị lộ cũng không chứa mã hợp lệ. Yêu cầu nhận 401 nếu phiên đã bị thu hồi, quá 12 giờ kể từ khi cấp hoặc quá 30 phút không hoạt động, và phiên bị thu hồi ngay nếu tài khoản bị khóa, ngừng hoạt động hoặc là Giám sát viên không còn khu vực. Khi người dùng còn thao tác, giao diện làm mới phiên; mỗi lần làm mới cấp mã mới với hạn 12 giờ riêng và thu hồi mã cũ, nên mã bị lộ chỉ có giá trị trong thời gian ngắn.

\subsection{Chống giả mạo yêu cầu}
\label{sec:5.5.2}

Vì trình duyệt tự gửi cookie, một trang web khác có thể khiến nó gửi yêu cầu tới hệ thống (CSRF). Ngoài \texttt{SameSite}, mọi yêu cầu thay đổi dữ liệu phải kèm tiêu đề \texttt{X-CSRF-Token}, là HMAC-SHA256 của mã phiên được trả cho giao diện khi đăng nhập, mà trang khác không đọc được. Chỉ các nguồn được cấu hình mới gọi được API kèm cookie (CORS), và phản hồi kèm các tiêu đề bảo mật cấm đoán kiểu nội dung, cấm nhúng trang và không lưu bộ nhớ đệm.

\subsection{Phân quyền}
\label{sec:5.5.3}

Phân quyền có ba lớp tại máy chủ. Mỗi API khai báo vai trò được phép và từ chối vai trò khác với mã 403 trước mọi xử lý. Tầng nghiệp vụ kiểm tra từng đối tượng: camera và lần xuất hiện thuộc khu vực hiện tại của Giám sát viên và đang vận hành, vụ việc do chính người đó phụ trách, ảnh theo một trong hai căn cứ ở Mục~\ref{sec:4.6}; khu vực và người phụ trách luôn lấy từ phiên chứ không từ yêu cầu, nên sửa yêu cầu ở trình duyệt không mở rộng được quyền. Cuối cùng, ràng buộc cơ sở dữ liệu chặn dữ liệu sai kể cả khi tầng nghiệp vụ có lỗi.

Đối tượng ngoài phạm vi nhận 404 thay vì 403, để việc thử lần lượt các mã không cho biết đối tượng nào tồn tại.

\subsection{Bảo vệ dữ liệu và hạn chế lạm dụng}
\label{sec:5.5.4}

Thông tin xác thực RTSP được tách khỏi địa chỉ, mã hóa bằng Fernet với khóa từ biến môi trường và chỉ giải mã khi kết nối; nó không bao giờ được trả về hay ghi nhật ký. Vì kiểm tra kết nối khiến máy chủ kết nối tới địa chỉ do người dùng nhập, chỉ địa chỉ IP thuộc dải mạng camera đã khai báo được chấp nhận, còn loopback, link-local và multicast luôn bị từ chối, để chức năng này không bị dùng dò mạng nội bộ (SSRF).

Các chức năng tốn tài nguyên hoặc dễ bị dò được giới hạn mỗi phút: đăng nhập 10 lần theo cặp IP - tên đăng nhập, và theo tài khoản: tìm kiếm 30, tải video 10, kiểm tra kết nối 10, kiểm tra AI 6 lần; vượt giới hạn thì nhận 429 kèm thời gian chờ. Bộ đếm nằm trong bộ nhớ tiến trình, phù hợp triển khai một tiến trình. Ảnh truy vấn phải là JPEG hoặc PNG, không quá 10~MB và 24 triệu điểm ảnh, video không quá 500~MB và có chữ ký định dạng hợp lệ, và tên tệp người dùng gửi không dùng làm đường dẫn lưu trữ.
